\section{并行化工作流与核心性能挑战}

\subsection{从“程序太慢”到可验证的并行实现}

当已有程序运行过慢时，需要找出耗时部分，把适合的部分放到图形处理器上，并持续优化直到性能满足目标。并行化应形成测量驱动的闭环，而不能从“直接重写全部代码”开始。

\subsubsection{建立正确的串行基线}

首先保留一个能够产生正确结果的基线版本，并记录代表性输入上的运行时间。基线提供两种参照：一是验证并行版本是否保持相同计算语义，二是定义加速比的分母。没有可靠基线，后续的“更快”可能来自减少计算、降低精度或产生错误结果。

\subsubsection{性能剖析并识别热点}

使用\term{性能剖析}{profiling}工具测量各函数、循环或阶段的时间占比，找出真正主导总运行时间的\term{热点}{hotspot}。如果某段代码只占总时间的 1\%，即使把它加速 100 倍，对整体性能的贡献也很有限；第~\ref{l1:sec:amdahl} 节将用阿姆达尔定律严格表达这一点。

剖析还应区分计算、内存访问、数据传输和同步等待。某一函数“耗时长”不一定表示算术量大，也可能表示它在等待数据或共享资源。

\subsubsection{判断原算法能否暴露可扩展并行性}

有些算法可以直接把独立迭代分配给不同线程；另一些算法包含长依赖链，原始形式不适合并行执行。此时需要采用或重新设计\term{可扩展算法}{scalable algorithm}。可扩展性意味着：问题规模和处理资源增加时，可同时利用的有用工作也能增加，而串行协调、通信与资源争用不会迅速吞噬收益。

因此，并行化可能只是实现层面的映射，也可能要求算法层面的替换。若现有算法不适合并行执行，程序员需要选择结构完全不同的算法。

\subsubsection{映射工作并优化数据布局}

把工作划分为大量线程后，还要决定数据如何存放、何时移动以及哪些线程访问哪些数据，其目标是优化数据排列以最大化局部性。\term{局部性}{locality}表示程序在较短时间内重复访问相近或相关数据的倾向。良好的数据布局可以减少高延迟访问、提高缓存或片上存储的复用，并让大量线程更有效地使用全局内存带宽。

\subsubsection{测量驱动的性能调优}

\term{性能调优}{performance tuning}是迭代过程：修改实现、测量、解释瓶颈，再决定下一项修改。每次优化必须同时检查正确性和性能，不能仅依据代码看起来“更并行”就断言加速成立。

完整工作流可概括为：
\begin{equation*}
\begin{aligned}
\text{正确基线}
&\rightarrow \text{性能剖析}
\rightarrow \text{热点与并行性分析}
\rightarrow \text{算法或实现重构},\\
&\rightarrow \text{数据布局优化}
\rightarrow \text{测量与验证}
\rightarrow \text{继续迭代}.
\end{aligned}
\end{equation*}

最后还要检查代码的可移植性、可扩展性和可维护性。一个依赖特定设备偶然参数、只在单一输入上有效的实现，可能在下一代硬件或更大数据集上失效。

\subsection{挑战一：大规模并行要求规则性}

图形处理器最有效的执行形态是许多线程在同一阶段执行相同或相近的操作，只是处理不同数据。这种要求称为\term{规则性}{regularity}。规则性包含至少三个层面：
\begin{itemize}
  \item 工作结构规则：线程执行路径相近，不出现少数线程承担大量额外步骤；
  \item 数据访问规则：线程以可预测方式访问数据，避免大量随机、冲突或低复用访问；
  \item 阶段推进规则：线程组能在相近时间完成当前阶段，减少等待与闲置。
\end{itemize}

规则性不是说所有线程必须处理完全相同的数据，而是说硬件可以用共享的控制与调度机制高效驱动大量相似工作。若每个线程都执行截然不同的控制流，图形处理器复制大量算术通道的优势就难以兑现。

\subsection{挑战二：负载均衡}

设一个并行阶段包含 $n$ 个任务，第 $i$ 个任务的执行时间为 $T_i$。如果下一阶段必须等所有任务完成，则该阶段的完成时间为
\begin{equation}
  T_{\mathrm{phase}}=\max_{1\le i\le n} T_i.
  \label{l1:eq:parallel-phase-time}
\end{equation}

式~\eqref{l1:eq:parallel-phase-time} 表明，最快线程提前完成并不能缩短阶段时间；最慢线程决定何时能够继续。这就是\term{负载均衡}{load balance}的核心。

\begin{figure}[H]
  \centering
  \includegraphics[width=0.91\textwidth]{load_balance.png}
  \caption{规则并行与负载不均衡}
\end{figure}

为了量化不均衡，可以定义一个解释性指标
\begin{equation}
  E_{\mathrm{load}}
  =\frac{\sum_{i=1}^{n}T_i}{n\max_i T_i},
  \qquad 0<E_{\mathrm{load}}\le 1.
  \label{l1:eq:load-efficiency}
\end{equation}
若所有任务耗时相同，则 $E_{\mathrm{load}}=1$；若少数任务远慢于其余任务，该比值会降低，表示大量并行资源在等待最慢任务时处于空闲状态。式~\eqref{l1:eq:load-efficiency} 是用于理解负载不均衡的解释性指标；性能也可以用其他指标量化。

负载不均衡可能来自输入数据差异、分支路径差异、任务划分粒度过粗或某些任务争用更多资源。解决方向是重新分割任务、细化工作粒度或选择更规则的算法，使各执行单元承担的有效工作接近。

\subsection{挑战三：全局内存带宽的理想值与实际值}

\term{全局内存带宽}{global memory bandwidth}表示设备内存与处理器之间单位时间能够传输的数据量。硬件规格给出的峰值带宽是在理想访问模式下的上限，应用实际获得的\term{有效带宽}{effective bandwidth}可写为
\begin{equation}
  B_{\mathrm{eff}}=\frac{D_{\mathrm{useful}}}{T},
  \label{l1:eq:effective-bandwidth}
\end{equation}
其中 $D_{\mathrm{useful}}$ 是完成目标计算所需的有效数据量，$T$ 是相应时间。若程序重复搬运数据、访问模式分散或不能充分并行发起内存请求，则 $B_{\mathrm{eff}}$ 可能远低于峰值。

水坝与吸管的类比可以区分“理想”与“现实”：设备可能具备很宽的总带宽，但单个程序若组织不当，只能利用其中很小一部分。提高有效带宽依赖正确的内存模型、数据复用和访问模式。

\subsection{挑战四：冲突访问导致串行化与延迟}

当大量线程同时访问某个受限或关键资源时，请求必须排队，原本并行的工作被迫部分\term{串行化}{serialization}。国际机场入境检查可以类比这种瓶颈：许多航班可以同时到达，但若所有旅客必须通过数量有限的检查窗口，检查阶段就成为窄口，前面的并行到达能力无法转化为端到端吞吐量。

\begin{figure}[H]
  \centering
  \includegraphics[width=0.92\textwidth]{bandwidth_serialization.png}
  \caption{全局内存带宽与关键资源争用}
\end{figure}

从系统角度看，争用会产生两类损失：
\begin{enumerate}
  \item 请求排队，单个访问的等待时间增加；
  \item 等待资源的线程不能推进，执行通道可能空闲，整体吞吐量下降。
\end{enumerate}

海量并行系统对串行化尤其敏感。若只有一个线程，某个小的串行步骤只是该线程的一部分；若有成千上万线程同时汇聚到同一关键资源，该步骤会把大量并行执行压缩为一条队列。同步、共享数据更新与内存访问设计都必须控制这种汇聚效应。

\subsection{四类挑战之间的耦合}

规则性、负载均衡、带宽和争用不是相互独立的问题。一个不规则算法可能同时造成：不同线程工作量差异、数据访问失去局部性、多个线程争用热点数据，以及难以隐藏的等待。相反，规则的数据并行结构通常更容易均匀划分工作，也更容易形成高效访存。

\begin{keybox}
并行化的目标不是把串行语句机械地复制到更多线程，而是构造一种硬件能够持续推进的工作形态：线程数量足够多，执行路径足够规则，任务负载相近，数据访问能够利用带宽，并且不存在把全部线程重新压回串行队列的关键资源。
\end{keybox}
